\subsection{The exact completed tower template}\label{sec:tower-template}
In one-based display notation define
\[
 H_{1,m}(q)=(\underbrace{q,\ldots,q}_{m}),\qquad
 H_{r+1,m}(q)=(q^{H_{r,m}(q+1)}),\qquad
 V_{r,m}=()(1^{H_{r,m}(2)}).
\]
Thus $V_{r,1}=I_{r+1}$. The completed root column has $m$ bare
entries followed by heads $G_1,\ldots,G_r$. The private column of
$G_j$ is a full intrinsic image of the root prefix consisting of
the $m$ bare entries and $G_1,\ldots,G_{j-1}$. The outer target column
is a full image of the whole root prefix. Every request pairs both
entry-by-entry, including its recursive private packages.

This is the original completion, not an added saturated table. Induct
on $j$. A request in the $j$th private package needs exactly the
equal-height same-parent roots in the earlier root prefix; they are
already present by the induction hypothesis. No strict candidate is
needed. Conversely each listed support has an actual reader in that
private package or in the outer image. All internal packages have the
same property recursively. Hence the table is completed and has the
least actual support. Lemma~\ref{lem:replay} identifies it with the
original completion of $V_{r,m}$.

Only the outer image of the last head is native at its level. Down its
private chain the last head remains native, ending in the $m$ native
bare entries. Other heads are virtual; source packages are globally
virtual but retain their nonempty relative birth templates. This
distinction explains why removing all globally virtual packages would
destroy the initialization. For reference, the completed finite shape
has $2^{r+1}$ columns and $2^{r+1}(m+1)-2$ entries, obtained by doubling
the preceding prefix and adding the next head. These are shape counts
only; the ordinal labels are supplied by the joint reflection in
Lemma~\ref{lem:prefix-init}.
